%=====================================================================
% Modelo de datos · Llaves primarias y restricciones de unicidad
%=====================================================================
\subsection*{Llaves primarias}
\addcontentsline{toc}{subsection}{Llaves primarias}

La llave primaria de cada entidad se eligió con cuatro criterios:

\begin{reglas}
  \item \textbf{Sustituta.} Un entero \at{IDENTITY} para las entidades con identidad propia, que se crean desde la aplicación y a las que otras tablas apuntan. Donde además existe un dato natural único —el \at{nickname}, el resumen de un token, la versión de términos de un usuario— se declara como llave candidata con \at{UNIQUE}.
  \item \textbf{Compuesta.} Las entidades asociativas y las débiles se identifican por las llaves que relacionan. La llave compuesta hace cumplir una regla sin código adicional: (\at{id\_usuario}, \at{id\_ranura}) en \textit{Equipamiento} garantiza un solo objeto por ranura (\mbox{CU-14} \mbox{RN-01}), y (\at{id\_usuario}, \at{id\_objeto}) en \textit{UsuarioObjeto} impide poseer dos veces el mismo objeto (\mbox{CU-38} \mbox{RN-02}).
  \item \textbf{Igual a la llave foránea.} Cuando solo puede haber una fila por usuario, la llave es el propio \at{id\_usuario}: un solo cambio de nombre (\mbox{CU-13} \mbox{RN-01}), un solo lugar en la cola y una sola sala a la vez (\mbox{CU-18} \mbox{RN-02}).
  \item \textbf{Asignada en la carga.} Los catálogos usan identificadores fijos, porque se cargan por SQL y la aplicación puede referirse a ellos por su valor.
\end{reglas}

% Archivo generado por bd/generar_tablas.ps1 a partir de bd/bastion.sql. No editar a mano.
\begin{center}\small
\begin{longtable}{|L{0.21\textwidth}|L{0.28\textwidth}|L{0.18\textwidth}|L{0.21\textwidth}|}
\hline
\textbf{Entidad} & \textbf{Llave primaria} & \textbf{Clase} & \textbf{Otras llaves candidatas} \\ \hline
\endfirsthead
\hline
\textbf{Entidad} & \textbf{Llave primaria} & \textbf{Clase} & \textbf{Otras llaves candidatas} \\ \hline
\endhead
\textit{Ranura} & (\texttt{id\_\allowbreak{}ranura}) & Simple, asignada en la carga & (\texttt{nombre}) \\ \hline
\textit{Objeto\-Cosmetico} & (\texttt{id\_\allowbreak{}objeto}) & Simple, asignada en la carga & --- \\ \hline
\textit{Modo} & (\texttt{id\_\allowbreak{}modo}) & Simple, asignada en la carga & (\texttt{nombre}) \\ \hline
\textit{Division} & (\texttt{id\_\allowbreak{}division}) & Simple, asignada en la carga & (\texttt{nombre}); (\texttt{orden}) \\ \hline
\textit{Nivel\-IA} & (\texttt{id\_\allowbreak{}nivel\_\allowbreak{}ia}) & Simple, asignada en la carga & (\texttt{nombre}) \\ \hline
\textit{Leccion} & (\texttt{id\_\allowbreak{}leccion}) & Simple, asignada en la carga & (\texttt{orden}) \\ \hline
\textit{Tipo\-Caja} & (\texttt{id\_\allowbreak{}tipo\_\allowbreak{}caja}) & Simple, asignada en la carga & (\texttt{nombre}) \\ \hline
\textit{Tipo\-Caja\-Objeto} & (\texttt{id\_\allowbreak{}tipo\_\allowbreak{}caja}, \texttt{id\_\allowbreak{}objeto}) & Compuesta & --- \\ \hline
\textit{Palabra\-Prohibida} & (\texttt{id\_\allowbreak{}palabra}) & Simple, sustituta & (\texttt{termino}, \texttt{idioma}, \texttt{ambito}) \\ \hline
\textit{Motivo\-Reporte} & (\texttt{id\_\allowbreak{}motivo}) & Simple, asignada en la carga & (\texttt{nombre}) \\ \hline
\textit{Usuario} & (\texttt{id\_\allowbreak{}usuario}) & Simple, sustituta & (\texttt{nickname}) \\ \hline
\textit{Sesion} & (\texttt{id\_\allowbreak{}sesion}) & Simple, sustituta & (\texttt{token\_\allowbreak{}hash}) \\ \hline
\textit{Codigo\-Segundo\-Factor} & (\texttt{id\_\allowbreak{}codigo}) & Simple, sustituta & --- \\ \hline
\textit{Token\-Verificacion\-Correo} & (\texttt{id\_\allowbreak{}token}) & Simple, sustituta & (\texttt{token\_\allowbreak{}hash}) \\ \hline
\textit{Token\-Recuperacion} & (\texttt{id\_\allowbreak{}token}) & Simple, sustituta & --- \\ \hline
\textit{Aceptacion\-Terminos} & (\texttt{id\_\allowbreak{}aceptacion}) & Simple, sustituta & (\texttt{id\_\allowbreak{}usuario}, \texttt{version\_\allowbreak{}terminos}) \\ \hline
\textit{Historial\-Nickname} & (\texttt{id\_\allowbreak{}usuario}) & Simple, igual a la llave foránea & --- \\ \hline
\textit{Avatar} & (\texttt{id\_\allowbreak{}avatar}) & Simple, sustituta & --- \\ \hline
\textit{Enlace\-Red} & (\texttt{id\_\allowbreak{}enlace}) & Simple, sustituta & (\texttt{id\_\allowbreak{}usuario}, \texttt{orden}) \\ \hline
\textit{Usuario\-Objeto} & (\texttt{id\_\allowbreak{}usuario}, \texttt{id\_\allowbreak{}objeto}) & Compuesta & --- \\ \hline
\textit{Equipamiento} & (\texttt{id\_\allowbreak{}usuario}, \texttt{id\_\allowbreak{}ranura}) & Compuesta & (\texttt{id\_\allowbreak{}usuario}, \texttt{id\_\allowbreak{}objeto}) \\ \hline
\textit{Partida} & (\texttt{id\_\allowbreak{}partida}) & Simple, sustituta & --- \\ \hline
\textit{Participacion} & (\texttt{id\_\allowbreak{}partida}, \texttt{id\_\allowbreak{}usuario}) & Compuesta & (\texttt{id\_\allowbreak{}partida}, \texttt{orden\_\allowbreak{}turno}) \\ \hline
\textit{Participacion\-Objeto} & (\texttt{id\_\allowbreak{}partida}, \texttt{id\_\allowbreak{}usuario}, \texttt{id\_\allowbreak{}ranura}) & Compuesta & (\texttt{id\_\allowbreak{}partida}, \texttt{id\_\allowbreak{}usuario}, \texttt{id\_\allowbreak{}objeto}) \\ \hline
\textit{Jugada} & (\texttt{id\_\allowbreak{}partida}, \texttt{numero\_\allowbreak{}jugada}) & Compuesta & --- \\ \hline
\textit{Oferta\-Tablas} & (\texttt{id\_\allowbreak{}oferta}) & Simple, sustituta & --- \\ \hline
\textit{Enlace\-Espectador} & (\texttt{id\_\allowbreak{}enlace}) & Simple, sustituta & (\texttt{token\_\allowbreak{}hash}) \\ \hline
\textit{Cola\-Emparejamiento} & (\texttt{id\_\allowbreak{}usuario}) & Simple, igual a la llave foránea & --- \\ \hline
\textit{Sala} & (\texttt{id\_\allowbreak{}sala}) & Simple, sustituta & --- \\ \hline
\textit{Sala\-Participante} & (\texttt{id\_\allowbreak{}usuario}) & Simple, igual a la llave foránea & (\texttt{id\_\allowbreak{}sala}, \texttt{plaza}) \\ \hline
\textit{Invitacion} & (\texttt{id\_\allowbreak{}invitacion}) & Simple, sustituta & --- \\ \hline
\textit{Mensaje} & (\texttt{id\_\allowbreak{}mensaje}) & Simple, sustituta & --- \\ \hline
\textit{Estadistica\-Modo} & (\texttt{id\_\allowbreak{}usuario}, \texttt{id\_\allowbreak{}modo}) & Compuesta & --- \\ \hline
\textit{Historial\-Division} & (\texttt{id\_\allowbreak{}historial\_\allowbreak{}division}) & Simple, sustituta & --- \\ \hline
\textit{Caja} & (\texttt{id\_\allowbreak{}caja}) & Simple, sustituta & --- \\ \hline
\textit{Caja\-Contenido} & (\texttt{id\_\allowbreak{}caja}, \texttt{numero}) & Compuesta & --- \\ \hline
\textit{Movimiento\-Moneda} & (\texttt{id\_\allowbreak{}movimiento}) & Simple, sustituta & --- \\ \hline
\textit{Amistad} & (\texttt{id\_\allowbreak{}usuario\_\allowbreak{}a}, \texttt{id\_\allowbreak{}usuario\_\allowbreak{}b}) & Compuesta & --- \\ \hline
\textit{Solicitud} & (\texttt{id\_\allowbreak{}solicitud}) & Simple, sustituta & (\texttt{id\_\allowbreak{}solicitante}, \texttt{id\_\allowbreak{}destinatario}) \\ \hline
\textit{Silencio} & (\texttt{id\_\allowbreak{}usuario}, \texttt{id\_\allowbreak{}silenciado}) & Compuesta & --- \\ \hline
\textit{Bloqueo} & (\texttt{id\_\allowbreak{}bloqueador}, \texttt{id\_\allowbreak{}bloqueado}) & Compuesta & --- \\ \hline
\textit{Progreso\-Tutorial} & (\texttt{id\_\allowbreak{}usuario}, \texttt{id\_\allowbreak{}leccion}) & Compuesta & --- \\ \hline
\textit{Reporte} & (\texttt{id\_\allowbreak{}reporte}) & Simple, sustituta & --- \\ \hline
\textit{Sancion} & (\texttt{id\_\allowbreak{}sancion}) & Simple, sustituta & --- \\ \hline
\textit{Apelacion} & (\texttt{id\_\allowbreak{}apelacion}) & Simple, sustituta & (\texttt{id\_\allowbreak{}sancion}) \\ \hline
\textit{Bitacora\-Moderacion} & (\texttt{id\_\allowbreak{}bitacora}) & Simple, sustituta & --- \\ \hline
\textit{Bitacora\-Acceso} & (\texttt{id\_\allowbreak{}bitacora}) & Simple, sustituta & --- \\ \hline
\end{longtable}
\end{center}


\subsubsection*{Restricciones de unicidad}
\addcontentsline{toc}{subsubsection}{Restricciones de unicidad}

Varias reglas de negocio dicen ``a lo sumo uno'' solo para las filas en cierto estado: un token vigente por cuenta, un avatar vigente, una participación sin terminar. No son llaves candidatas, porque las filas antiguas se conservan, así que se declaran como índices únicos filtrados: la unicidad se exige únicamente entre las filas que cumplen la condición. La tabla reúne todas las restricciones de unicidad del esquema, con la regla que cada una hace cumplir.

% Archivo generado por bd/generar_tablas.ps1 a partir de bd/bastion.sql. No editar a mano.
\begin{center}\small
\begin{longtable}{|L{0.2\textwidth}|L{0.25\textwidth}|L{0.45\textwidth}|}
\hline
\textbf{Entidad} & \textbf{Columnas y condición} & \textbf{Regla que garantiza} \\ \hline
\endfirsthead
\hline
\textbf{Entidad} & \textbf{Columnas y condición} & \textbf{Regla que garantiza} \\ \hline
\endhead
\textit{Ranura} & (\texttt{nombre}) & Dos ranuras no se llaman igual. \\ \hline
\textit{Objeto\-Cosmetico} & (\texttt{id\_\allowbreak{}objeto}, \texttt{id\_\allowbreak{}ranura}) & Superclave que usan Equipamiento y ParticipacionObjeto para exigir que el objeto sea de la ranura. \\ \hline
\textit{Objeto\-Cosmetico} & (\texttt{id\_\allowbreak{}ranura}), si \texttt{es\_\allowbreak{}predeterminado = 1} & Un solo objeto predeterminado por ranura (\mbox{CU-02} \mbox{PRE-05}). \\ \hline
\textit{Modo} & (\texttt{nombre}) & Dos modos no se llaman igual. \\ \hline
\textit{Division} & (\texttt{nombre}) & Dos divisiones no se llaman igual. \\ \hline
\textit{Division} & (\texttt{orden}) & Dos divisiones no ocupan la misma posición. \\ \hline
\textit{Nivel\-IA} & (\texttt{nombre}) & Dos niveles no se llaman igual. \\ \hline
\textit{Leccion} & (\texttt{orden}) & Dos lecciones no ocupan la misma posición. \\ \hline
\textit{Tipo\-Caja} & (\texttt{nombre}) & Dos tipos de caja no se llaman igual. \\ \hline
\textit{Palabra\-Prohibida} & (\texttt{termino}, \texttt{idioma}, \texttt{ambito}) & Un término no se repite en el mismo idioma y ámbito. \\ \hline
\textit{Motivo\-Reporte} & (\texttt{nombre}) & Dos motivos no se llaman igual. \\ \hline
\textit{Usuario} & (\texttt{nickname}) & El nickname es identificador de acceso (\mbox{CU-01} \mbox{RN-14}). \\ \hline
\textit{Usuario} & (\texttt{correo}), si \texttt{correo IS NOT NULL} & El correo es único entre las cuentas que lo tienen (\mbox{CU-02} \mbox{RN-02}). \\ \hline
\textit{Usuario} & (\texttt{codigo\_\allowbreak{}amigo}), si \texttt{codigo\_\allowbreak{}amigo IS NOT NULL} & El código de amigo es único por cuenta (\mbox{CU-30} \mbox{RN-06}). \\ \hline
\textit{Sesion} & (\texttt{token\_\allowbreak{}hash}) & Un token identifica una sola sesión. \\ \hline
\textit{Codigo\-Segundo\-Factor} & (\texttt{id\_\allowbreak{}usuario}), si \texttt{estado = PENDIENTE} & Un solo código vigente por cuenta (\mbox{CU-01} \mbox{RN-09}). \\ \hline
\textit{Token\-Verificacion\-Correo} & (\texttt{token\_\allowbreak{}hash}) & El enlace identifica un solo token. \\ \hline
\textit{Token\-Verificacion\-Correo} & (\texttt{id\_\allowbreak{}usuario}), si \texttt{estado = PENDIENTE} & Un solo token de verificación vigente por cuenta (\mbox{CU-04} \mbox{RN-03}). \\ \hline
\textit{Token\-Recuperacion} & (\texttt{id\_\allowbreak{}usuario}), si \texttt{estado = PENDIENTE} & Un solo código de recuperación vigente por cuenta (\mbox{CU-03} \mbox{RN-05}). \\ \hline
\textit{Aceptacion\-Terminos} & (\texttt{id\_\allowbreak{}usuario}, \texttt{version\_\allowbreak{}terminos}) & Una cuenta acepta cada versión una sola vez. \\ \hline
\textit{Avatar} & (\texttt{id\_\allowbreak{}usuario}), si \texttt{vigente = 1} & A lo sumo un avatar vigente por cuenta (\mbox{CU-12} \mbox{RN-06}). \\ \hline
\textit{Enlace\-Red} & (\texttt{id\_\allowbreak{}usuario}, \texttt{orden}) & Junto con el CHECK del orden, limita a cuatro enlaces por cuenta (\mbox{CU-12} \mbox{RN-03}). \\ \hline
\textit{Equipamiento} & (\texttt{id\_\allowbreak{}usuario}, \texttt{id\_\allowbreak{}objeto}) & Llave candidata: como el objeto determina su ranura, la cuenta y el objeto también identifican la fila. \\ \hline
\textit{Participacion} & (\texttt{id\_\allowbreak{}partida}, \texttt{orden\_\allowbreak{}turno}) & Dos jugadores no comparten turno en la misma partida. \\ \hline
\textit{Participacion} & (\texttt{id\_\allowbreak{}usuario}), si \texttt{terminada = 0} & Una sola participación sin terminar por jugador, incluidas las de la IA (\mbox{D-11}). \\ \hline
\textit{Participacion\-Objeto} & (\texttt{id\_\allowbreak{}partida}, \texttt{id\_\allowbreak{}usuario}, \texttt{id\_\allowbreak{}objeto}) & Llave candidata: como el objeto determina su ranura, también identifica la fila. \\ \hline
\textit{Oferta\-Tablas} & (\texttt{id\_\allowbreak{}partida}), si \texttt{estado = PENDIENTE} & A lo sumo una oferta pendiente por partida; la segunda es aceptación (\mbox{CU-22} \mbox{RN-04}). \\ \hline
\textit{Enlace\-Espectador} & (\texttt{token\_\allowbreak{}hash}) & El token identifica un solo enlace. \\ \hline
\textit{Sala} & (\texttt{codigo}), si \texttt{estado <> CERRADA} & El código es único entre las salas no cerradas y puede reutilizarse después (\mbox{CU-18} \mbox{RN-01}). \\ \hline
\textit{Sala} & (\texttt{id\_\allowbreak{}partida}), si \texttt{id\_\allowbreak{}partida IS NOT NULL} & Una partida viene de una sola sala. \\ \hline
\textit{Sala\-Participante} & (\texttt{id\_\allowbreak{}sala}, \texttt{plaza}) & Dos jugadores no ocupan la misma plaza (\mbox{CU-19} \mbox{RN-02}). \\ \hline
\textit{Invitacion} & (\texttt{id\_\allowbreak{}emisor}, \texttt{id\_\allowbreak{}destinatario}), si \texttt{estado = PENDIENTE} & Una sola invitación pendiente del mismo emisor al mismo destinatario (\mbox{CU-32} \mbox{RN-04}). \\ \hline
\textit{Solicitud} & (\texttt{id\_\allowbreak{}solicitante}, \texttt{id\_\allowbreak{}destinatario}) & No hay dos solicitudes iguales entre el mismo par (\mbox{CU-30} \mbox{RN-03}). \\ \hline
\textit{Reporte} & (\texttt{id\_\allowbreak{}denunciante}, \texttt{id\_\allowbreak{}reportado}, \texttt{id\_\allowbreak{}partida}), si \texttt{id\_\allowbreak{}partida IS NOT NULL} & No se reporta dos veces al mismo jugador por la misma partida (\mbox{CU-42} \mbox{RN-03}). \\ \hline
\textit{Apelacion} & (\texttt{id\_\allowbreak{}sancion}) & Una sola apelación por sanción (\mbox{CU-45} \mbox{RN-01}). \\ \hline
\end{longtable}
\end{center}

